\chapter{Contexte Général du Projet}

\section*{Introduction du chapitre}
Ce premier chapitre a pour vocation de planter le décor du projet \textbf{QuickTriage}.
Il s'agit de comprendre l'environnement dans lequel s'inscrit la solution, à savoir le
fonctionnement des services d'urgences hospitalières et les difficultés structurelles
auxquelles ils font face. Après avoir exposé le contexte et la problématique, nous
présenterons la motivation qui a guidé ce travail, puis nous formaliserons les objectifs,
principal et spécifiques. Le chapitre se clôt par la mise en évidence des contributions et
des innovations apportées par QuickTriage.

\section{Contexte et Problématique}

\subsection{Le triage aux urgences : une fonction critique}
Le service des urgences est, par définition, un environnement où l'incertitude et le
caractère vital des décisions cohabitent en permanence. À leur arrivée, les patients
présentent des tableaux cliniques extrêmement hétérogènes, allant du motif bénin à
l'urgence vitale. Le \textit{triage} consiste à évaluer rapidement la gravité de chaque
patient afin de déterminer l'ordre et le délai de prise en charge. Cette étape, réalisée
généralement par un infirmier organisateur de l'accueil (IOA), conditionne l'ensemble du
parcours de soins : un triage juste sauve des vies, un triage erroné peut être fatal.

\subsection{L'engorgement des urgences}
Partout dans le monde, les services d'urgences connaissent une croissance soutenue du
nombre de passages, sans augmentation proportionnelle des moyens. Cet
\textit{engorgement} se traduit par des temps d'attente allongés, une saturation des box
de soins et une charge mentale accrue pour le personnel. Dans un tel contexte, la pression
temporelle pesant sur la décision de triage augmente, et avec elle le risque d'erreur.
La rapidité et la fiabilité de l'évaluation initiale deviennent des enjeux de premier plan.

\subsection{Les limites du triage manuel}
Le triage manuel, bien qu'encadré par des échelles standardisées, demeure tributaire de
plusieurs facteurs de variabilité :
\begin{itemize}
  \item \textbf{La subjectivité de l'évaluateur} : deux soignants peuvent attribuer des
  niveaux différents à un même patient, en fonction de leur expérience et de leur état de
  fatigue.
  \item \textbf{La pression temporelle} : la nécessité de décider vite favorise les biais
  cognitifs et les raccourcis de raisonnement.
  \item \textbf{La charge cognitive} : l'intégration mentale simultanée de nombreux
  paramètres vitaux est difficile, surtout en période d'affluence.
  \item \textbf{Le risque de sous-triage} : sous-estimer la gravité d'un patient constitue
  l'erreur la plus dangereuse, car elle retarde une prise en charge potentiellement vitale.
\end{itemize}

\begin{warningbox}
Le \textit{sous-triage} (classer un patient grave comme non urgent) est l'erreur la plus
redoutée en médecine d'urgence. La conception de QuickTriage place donc la maximisation du
\textbf{rappel sur la classe \textsc{Critique}} au cœur de ses priorités, quitte à tolérer
quelques sur-triages.
\end{warningbox}

\subsection{Énoncé de la problématique}
La problématique centrale de ce projet peut être formulée ainsi : \textit{Comment concevoir
un outil numérique fiable, rapide et accessible, capable d'assister le personnel soignant
dans la décision de triage en objectivant l'évaluation de la priorité d'un patient à partir
de ses paramètres vitaux, tout en garantissant que la décision clinique finale demeure de la
responsabilité du soignant~?}

\section{Motivation du Projet}
La motivation du projet QuickTriage procède d'un triple constat. D'abord, un
\textbf{constat clinique} : les erreurs de triage, et singulièrement le sous-triage,
demeurent une cause évitable de morbidité. Ensuite, un \textbf{constat technologique} :
les progrès récents de l'apprentissage automatique, conjugués à des formats d'échange de
modèles interopérables comme ONNX et à des frameworks mobiles modernes comme Flutter,
rendent désormais réaliste le déploiement d'un outil d'aide à la décision léger et
embarquable. Enfin, un \textbf{constat d'opportunité} : il existe des jeux de données
publics de triage (notamment basés sur l'échelle KTAS) permettant d'entraîner et de valider
de tels modèles.

L'ambition n'est pas de remplacer le soignant, mais de lui offrir un \textbf{second avis
objectif et instantané}, agissant comme un filet de sécurité contre les erreurs
d'inattention et comme un accélérateur de la décision en période de forte affluence.

\section{Objectifs}

\subsection{Objectif principal}
L'objectif principal est de \textbf{concevoir et réaliser un système complet d'aide au
triage} qui prédit automatiquement le niveau de priorité KTAS d'un patient à partir de ses
constantes vitales, et le restitue de manière claire et immédiate au personnel soignant via
une application mobile.

\subsection{Objectifs spécifiques}
\begin{enumerate}
  \item \textbf{Construire un modèle d'apprentissage automatique} robuste de classification
  du niveau KTAS, en traitant explicitement les difficultés de qualité des données (fuite de
  données, déséquilibre des classes).
  \item \textbf{Maximiser le rappel de la classe \textsc{Critique}} afin de minimiser le
  sous-triage, par une calibration adaptée du seuil de décision.
  \item \textbf{Exposer le modèle via une API REST} performante et interopérable, en
  s'appuyant sur le format ONNX et une architecture J2EE (Jakarta EE / WildFly).
  \item \textbf{Développer une application mobile Flutter} ergonomique, optimisée pour une
  saisie rapide et une lecture immédiate du résultat (code couleur de priorité).
  \item \textbf{Garantir la primauté du jugement clinique}, le système se positionnant
  strictement comme une aide à la décision et non comme un décideur autonome.
\end{enumerate}

\section{Contributions et Innovations}
Au-delà de l'assemblage de briques technologiques, le projet QuickTriage se distingue par
plusieurs contributions :
\begin{itemize}
  \item \textbf{Diagnostic et correction d'une fuite de données.} Une première version du
  modèle s'appuyait sur un jeu de données contenant une variable (\texttt{plasma\_glucose})
  introduisant une fuite de données (\textit{data leakage}). Sa détection a conduit à un
  changement de jeu de données vers \textit{KTAS Emergency Triage}, gage de résultats
  honnêtes et généralisables.
  \item \textbf{Stratégie d'augmentation hybride des données.} Pour pallier le déséquilibre
  des classes, une approche combinant la génération de patients synthétiques calibrés sur la
  distribution réelle et l'algorithme SMOTE (avec \texttt{k\_neighbors} adaptatif) a été mise
  en œuvre, aboutissant à un jeu équilibré de 3~000 patients (600 par classe).
  \item \textbf{Calibration orientée sécurité patient.} Le seuil de décision de la classe
  \textsc{Critique} a été abaissé à 0,30 afin de privilégier le rappel des cas graves, en
  cohérence avec l'impératif clinique de limiter le sous-triage.
  \item \textbf{Chaîne d'intégration interopérable.} L'export du modèle au format ONNX et son
  chargement unique au démarrage du serveur J2EE assurent des inférences rapides et
  reproductibles, indépendamment du langage d'entraînement.
\end{itemize}

\begin{successbox}
QuickTriage n'est pas un simple prototype de modèle~: c'est une \textbf{chaîne complète et
cohérente}, de la donnée brute jusqu'à l'application mobile, dont chaque maillon a été
pensé pour la fiabilité et la sécurité du patient.
\end{successbox}

\section*{Conclusion du chapitre}
Ce chapitre a permis de cerner le contexte du triage aux urgences, d'en exposer les limites
actuelles et de formuler la problématique à laquelle QuickTriage entend répondre. Nous avons
également défini les objectifs du projet et mis en lumière ses principales contributions.
Pour situer notre solution par rapport à l'existant et en justifier les choix, le chapitre
suivant propose un état de l'art des systèmes de triage et un positionnement comparatif de
QuickTriage.
